% \documentclass{12pt,a4paper}{article}
% \usepackage{utf8}{inputenc}
% \usepackage{ngerman}{babel}
% \usepackage{xcolor}
% \usepackage{setspace}
% \usepackage{margin=2.5cm}{geometry}

\author{OTS}
\documentclass[12pt]{../../inc/mybib}

\setincpath{../../inc/}

\usepackage{bible_style}
\usepackage{textanalysis}

\graphicspath{{../../assets/images/}}

% Absatzformatierung
\setlength{\parindent}{0pt}
\setlength{\parskip}{1em}

\begin{document}

\section{Textanalysen}
\subsection{Hebräer 4,12-16}
\subsubsection{Text}
\textbf{12} \konj{Denn}\annot{Ursache} das Wort \pers{Gottes} \verbstamm{ist} \annot{Präsens: Indikativ} lebendig 

    \hspace*{1cm}\konj{Und}\annot{Aufreihung} wirksam 

    \hspace*{1cm}\konj{Und}\annot{Aufreihung} schärfer als jedes zweischneidige Schwert, 

        \hspace*{2cm}\konj{und}\annot{Vortsetzung} es \annot{Das Wort} \verbstamm{dringt durch}\annot{Präsens: Indikativ}, 

            \hspace*{3cm}\konj{bis}\annot{Zeit} es\annot{Das Wort} \verbstamm{scheidet}\annot{Futur: Konjunktiv}, 

                \hspace*{4cm}\konj{sowohl}\annot{Nebenordnend} Seele als auch Geist,

                \hspace*{4cm}\konj{sowohl}\annot{Nebenordnend} Mark als auch Bein,

        \hspace*{2cm}\konj{und}\annot{Vortsetzung} es\annot{Das Wort} \verbstamm{ist}\annot{Präsens: Indikativ} ein \pers{Richter} der Gedanken 

        \hspace*{2cm}\konj{und}\annot{Aufreihung} Gesinnungen des Herzens.

\textbf{13} \konj{Und}\annot{Vortsetzung} kein \ort{Geschöpf} \verbstamm{ist}\annot{Präsens: Indikativ} vor ihm\annot{Das Wort} \verbstamm{verborgen}\annot{Partizip 2}, 

    \hspace{1cm}\konj{sondern}\annot{Folge} alles \verbstamm{ist}\annot{Präsens: Indikativ} \verbstamm{enthüllt}\annot{Partizip 2}

        \hspace{2cm}\konj{und}\annot{Aufzählung} \verbstamm{aufgedeckt}\annot{Partizip 2} vor den Augen dessen\annot{WESSEN AUGEN?}, 

        \hspace{2cm}dem\annot{dessen} wir Rechenschaft zu \verbstamm{geben}\annot{Infinitiv} \verbstamm{haben}\annot{Präsens: Indikativ}.

\textbf{14} \konj{Da}\annot{Quelle} wir \konj{nun}\annot{ZEIT} einen Grossen \pers{Hohepriester} \verbstamm{haben}\annot{Präsens: Indikativ},

    \hspace{1cm}der\pers{Jesus}\annot{Hohepriester} die \ort{Himmel} \verbstamm{durchschritten hat}\annot{Perfekt: Indikativ}, \pers{Jesus, den Sohn Gottes}, 

    \hspace{1cm}\konj{so}\annot{Folge} \verbstamm{lasst}\annot{Präsens: Imperativ} uns \verbstamm{festhalten}\annot{Infinitiv} an dem Bekenntnis!

\textbf{15} \konj{Denn}\annot{Begründung} wir \verbstamm{haben}\annot{Präsens: Indikativ} nicht einen \pers{Hohepriester}\annot{Jesus},

    \hspace{1cm}der kein Mitleid \verbstamm{haben}\annot{Infinitv} \verbstamm{könnte}\annot{Präteritum} mit unseren Schwachheiten, 

        \hspace{2cm}\konj{sondern}\annot{Gegensatz} einen\pers{Hohepriester}, 

            \hspace{3cm}der\pers{Hohepriester},  

            \hspace{3cm}in allem \verbstamm{versucht worden ist}\annot{Passiv}

            \hspace{3cm}in ähnlicher Weise\annot{wie wir}, 

doch{Einschränkung} ohne Sünde.

\textbf{16} \konj{So}{Folge} \verbstamm{lasst}\annot{Präsens: Imperativ} uns 
Nun{Zeit} mit Freimütigkeit hinzutreten{Partizip 2} zum Thron der Gnade. 
\konj{Damit}\annot{Ergebnis} wir Barmherzigkeit erlangen{Präsens: Infinitiv} 
\konj{Und}{Aufzählung} Gnade \verbstamm{finden} zu rechtzeitiger Hilfe!

\subsubsection{Hinweise}
\begin{itemize}     
    \item Wort Gottes ist: V12
    \begin{itemize}       
        \item Wirksam
        \item Scharf wie ein Schwert
        \item Richter
        \item Deckt auf (Gesinnung)
        \item Niemand ist vor ihn Verborgen V13
    \end{itemize}
    \item Wir haben eine Grossen Hohepriester V14
    \item Wir können Freimütig sein
\end{itemize}

\subsubsection{Beobachtungen}
\begin{itemize}
    \item Lebendiges Wort und das Wort kommt von Gott.
    \item Das Wort ist wirksam
    \item Das Wort ist ein Spiegel für unsere Seele und für unser Leben.
    \item Das Wort durchdringt unser ganzes Wesen.
    \item Das Wort ist ein Richter wir können uns nach diesem Wort Gottes ausrichten.
    \item Alle Geschöpfe, alles was Geschaffen ist, gehört dem der das Wort geschaffen hat.
    \item Wir haben einen Hohepriester, einen Führsprecher der bis ins Allerheiligste eingedrungen ist.
    \item Dieser Hohepriester ist Jesus Christus.
    \item Dieser Fürsprecher kennt uns. Jesus hat das gleiche durchgemacht wie wir als er als Mensch auf die \item Erde kam. Ja noch viel mehr.
\item Wegen diesem Führsprecher, können wir getrost vor Gott erscheinen.
\end{itemize}

\subsubsection{Aussage / Absicht}
\renewcommand{\arraystretch}{1.8}
\begin{longtable}{l p{10cm} p{10cm}}
    \hline
    \textbf{Vers} &	\textbf{Aussage} & \textbf{Absicht} \\
    \hline
    \endfirsthead
    \hline
    \textbf{Vers} &	\textbf{Aussage} & \textbf{Absicht} \\
    \hline
    \endhead
    12 & Denn das Wort Gottes ist lebendig und wirksam und schärfer als jedes zweischneidige Schwert, und es dringt durch, bis es scheidet \konj{sowohl} Seele als auch Geist, \konj{sowohl} Mark als auch Bein, und es ist ein Richter der Gedanken und Gesinnung des Herzen.
    & Wir sollen das Wort Gottes studieren, weil es uns unsere Gesinnung und Gedanken prüft. Das Wort kommt von Gott. \\
    \hline
    13 &	Niemand ist vor dem Wort Gottes verborgen &	Das sich jeder Bewusst ist, dass alle Geschöpfe sich an diesem Wort messen müssen. Eines Tages werden wir nach diesem Wort Gerichtet. \\
    \hline
    14	& Jesus Christus unser Hohepriester	& Für die Juden war der Hohepriester der, der 1x im Jahr ins Heiligtum ging um für die Sünden von ganz Israel ein Opfer zu bringen. \\
    \hline
    14	& Festhalten &	Wir sollen uns an diesem Hoheprister, der die Himmel durchschritten hat, also gestorben und auferstanden ist festhalten. \\
    \hline
    15 & Der Hohepriester war Mensch und einer wie wir, nur ohne Sünde. &	Dadurch das Jesus ein Mensch war, versteht er uns und unsere Nöte. Er hat alles durchlebt und durchlitten. Aber er ist vollkommen und absolut vertrauenswürdig. \\
    \hline
    16	& Freimütig und Frohen Herzens &	Wir brauchen keine Angst zu haben wenn wir uns an Jesus angehängt haben. Er ist unser Führsprecher, wenn wir nach dem Wort gerichtet werden.\\
    \hline
\end{longtable}


\subsection*{Josua 3,1-5}

\textbf{1} \konj{Da} \verbstamm{machte} \annot{Präteritum-Indikativ-3P-Aktiv-Singular} \nomen{Josua} \partikel{früh} \verbpart{auf} \annot{Verbpartikel},

\hspace*{1cm} \konj{und} \plural{sie} \annot{das-Volk} \verbstamm{zogen} \annot{Präteritum-Indikativ-3P-Plural} \prep{aus} \annot{Verbpartikel} \plural{Sittim}, ¶

\hspace*{1cm} \konj{und} \verbstamm{kamen} \annot{Präteritum-Indikativ, Aktiv, 3P-Plural} \prep{an} \prep{den} \nomen{Jordan}, ¶

\hspace*{2cm} er \annot{Josua} \konj{und} \plural{alle} \nomen{Kinder} \nomen{Israels} \annot{das-Volk}; ¶

\hspace*{2cm} \konj{und} \plural{sie} \annot{Volk-und-Josua} \verbstamm{rasteten} \annot{Präteritum, Aktiv, Indikativ, 3p, Plural} \ort{dort}, \konj{ehe} \plural{sie} \annot{Volk-und-Josua} \verbstamm{hinüberzogen}. ¶

¶

\textbf{2} \konj{Nach} \prep{drei} \nomen{Tage} \konj{aber} ¶

\hspace*{1cm} \verbstamm{gingen} \prep{die} \nomen{Vorsteher} \prep{durch} \prep{das} \ort{Lager} ¶

\hspace*{2cm} \textbf{3} \konj{und} \verbstamm{geboten} \prep{dem} \nomen{Volk}: ¶

\hspace*{2cm} \konj{und} \verbstamm{sprachen}: ¶

\hspace*{3cm} \textbf{Wenn} \prep{ihr} \prep{die} \nomen{Bundeslade} \prep{des} \nomen{Herrn}, \partikel{eures} \nomen{Gottes}, \verbstamm{sehen} \verbstamm{werdet} ¶

\hspace*{4cm} \konj{und} \prep{die} \nomen{Priester}, \prep{die} \nomen{Leviten}, \prep{die} \plural{sie} \verbstamm{tragen}, ¶

\hspace*{4cm} \konj{so} \verbstamm{brecht} \verbpart{auf} \prep{von} \partikel{eurem} \nomen{Ort} ¶

\hspace*{5cm} \konj{und} \nomen{folgt} \prep{ihr} \verbpart{nach}! ¶

\hspace*{3cm} \textbf{4} \konj{Doch} \konj{soll} \prep{zwischen} \partikel{euch} ¶

\hspace*{3cm} \konj{und} \prep{ihr} \partikel{etwa} \nomen{2000} \nomen{Ellen} \nomen{Abstand} \verbstamm{sein}. ¶

\hspace*{2cm} \nomen{Kommt} \prep{ihr} \konj{nicht} \prep{zu} \ort{nahe}, ¶

\hspace*{2cm} \konj{damit} \prep{ihr} \prep{den} \nomen{Weg} \verbstamm{erkennt}, \prep{den} \prep{ihr} \verbstamm{gehen} \konj{sollt}; ¶

\hspace*{3cm} \konj{denn} \prep{ihr} \verbstamm{seid} \prep{den} \nomen{Weg} \konj{zuvor} \verbstamm{gegangen}! ¶

¶

\textbf{5} \konj{Und} \nomen{Josua} \verbstamm{sprach} \prep{zum} \nomen{Volk}: \nomen{Heiligt} \partikel{euch}, ¶

\vspace{2em}

\subsection*{Legende}

\begin{itemize}
    \item \verbstamm{Grün} -- Verbstämme und Verbformen
    \item \nomen{Orange} -- Nomen und Substantive
    \item \konj{Rot} -- Konjunktionen und Bindewörter
    \item \prep{Blau} -- Präpositionen, Artikel und Partikeln
    \item \plural{Magenta} -- Pluralformen und Pronomen
    \item \ort{Lila} -- Ortsangaben
    \item \annot{Grau, kursiv} -- Grammatische Annotationen
\end{itemize}

\vspace{1em}
\textit{¶ = Textsegmentierung/Strukturmarkierung}

\end{document}
